\documentclass[10pt,twocolumn,letterpaper]{article}

\usepackage[utf8]{inputenc}
\usepackage[T1]{fontenc}
\usepackage{lmodern}
\usepackage{amsmath,amssymb}
\usepackage{graphicx}
\usepackage{booktabs}
\usepackage[margin=1in]{geometry}
\usepackage{authblk}
\usepackage[colorlinks=true, allcolors=blue]{hyperref}
\usepackage{float}
\usepackage{cite}
\usepackage{caption}
\usepackage{placeins}

\title{\Large \textbf{In Silico Design and Computational Validation of a Tri-Modular Genetic Circuit for Autonomous Detection and Bioremediation of Hexavalent Chromium in Closed-System Microcosms}}

\author{Zubayer Hasan Shaad\thanks{Corresponding author: Zubayer Hasan Shaad, Govt. Tolaram College, Narayanganj, Bangladesh. Email: mdzubayerhasanshaad99@gmail.com}}
\affil{Govt. Tolaram College, Narayanganj, Bangladesh}
\date{July 2026}

\begin{document}

\twocolumn[
  \maketitle
  \begin{center}
  \large\textbf{Abstract}
  \end{center}
  \noindent
  Hexavalent chromium [Cr(VI)], a Group 1 IARC carcinogen, has been measured at concentrations up to 3.54 mg/L in Bangladesh's Buriganga River, approximately 35-fold above the US EPA surface water limit, primarily as a consequence of unregulated leather tannery effluent discharge. While existing remediation technologies address contamination passively, none integrate real-time sensing, dose-responsive enzymatic treatment, and autonomous post-treatment self-termination. More critically, prior proposals deploying engineered microorganisms in open rivers carry catastrophic ecological risks that remain computationally unquantified.

  \noindent
  Here, we present a fully computational design and validation of a tri-modular synthetic genetic circuit expressed in a double-auxotrophic \textit{Escherichia coli} chassis, exclusively deployed within a closed-system microcosm (bioreactor or sealed pond). The three modules, each encoded on an orthogonal compatibility-verified plasmid backbone, function as follows: \textbf{(1)} a ChrB-sfGFP biosensor module (pChrB-sfGFP, pUC19) provides Cr(VI) detection at nanomolar sensitivity via a de-repression logic gate; \textbf{(2)} a NemA chromate reductase module (pNemA-His, pET-28a), driven by the identical PchrB promoter, ensures \textit{coordinately regulated}, dose-responsive enzymatic Cr(VI) $\rightarrow$ Cr(III) reduction; \textbf{(3)} a dual-trigger holin-endolysin 'Deadman' kill switch (pKillSwitch-DT, pACYC184) enforces programmed cell death via AND-gate logic integrating chromate depletion sensing and a constitutive CI434-SsrA timer repressor with a 24-hour protein half-life.

  \noindent
  The computational validation framework spans four orthogonal methodologies: Ordinary Differential Equation (ODE) systems modeling of the gene expression dynamics across all three modules over 96 hours; comparative Michaelis-Menten modeling of a hypothesized NemA$^{*2+}$ catalytic variant whose kinetic parameters ($K_m^* = 16\ \mu$M, $k_{cat}^* = 37.2\ \text{h}^{-1}$) are \textbf{assumed} by analogy to OYE-family active-site mutagenesis precedent \cite{mowafy2010}, not derived from a structural modeling pipeline executed in this study; a ribosome-allocation genome-scale metabolic model demonstrating that the circuit's $\sim$4.5\% proteome burden is metabolically sustainable at environmentally relevant chromium concentrations; and Luria-Delbr\"uck fluctuation analysis showing that, under the assumed input parameters, the layered containment architecture (genetic kill switch + double auxotrophy $\Delta$thyA $\Delta$dapA) yields a 30-day probability of a viable escapee in a 1,000-liter closed system of $P_{\text{escape}} \approx 6 \times 10^{-17}$, below an acceptable-risk threshold of $10^{-15}$. A sensitivity sweep (Pillar 4b) shows this headline figure is robust to a 10$\times$ error in the assumed per-bp mutation rate and to a 10$\times$ scale-up of the deployment reactor, but is not robust to a 100$\times$ error in $\mu_{bp}$, at which point $P_{\text{escape}}$ rises to $6 \times 10^{-15}$, above the threshold. All simulations were executed in Python (SciPy, NumPy, Matplotlib).

  \noindent
  This work demonstrates that, under its modeling assumptions, the proposed tri-modular system can reduce 100 $\mu$M Cr(VI) by $>$99\% within 48 hours under closed-system conditions. We argue that this computational framework, which mathematically quantifies both remediation performance and mutation-driven failure modes, and which explicitly reports the sensitivity of its safety claim to its key input assumption, establishes a transparent in silico biosafety validation template for BMO-based heavy metal remediation in resource-constrained settings. The work is entirely computational; no wet-lab validation has been performed.

  \vspace{0.5cm}
  \begin{center}
  \textbf{Keywords:} synthetic biology $\cdot$ hexavalent chromium $\cdot$ genetic kill switch $\cdot$ ChrB biosensor $\cdot$ NemA reductase $\cdot$ bioremediation $\cdot$ Luria-Delbr\"uck mutation model $\cdot$ plasmid architecture $\cdot$ BMO biosafety $\cdot$ Bangladesh tannery pollution
  \end{center}
]

\section{Introduction}

\subsection{The Hexavalent Chromium Crisis in Bangladesh}
Hexavalent chromium [Cr(VI)] is a highly soluble, mobile, and persistent environmental contaminant classified by the International Agency for Research on Cancer (IARC) as a definitive Group 1 human carcinogen \cite{iarc1990}. At the molecular level, Cr(VI) enters cells through nonspecific anion transport channels (sharing uptake pathways with sulfate and phosphate), where it undergoes intracellular reduction to Cr(III) and Cr(V) intermediates. These intermediates generate reactive oxygen species (ROS) that form DNA adducts, strand breaks, and protein-DNA crosslinks, mechanisms directly linked to lung carcinoma, nasal septal perforation, and renal tubular necrosis upon chronic exposure \cite{eastmond2008}.

In Bangladesh, the leather tanning industry, historically centred in Hazaribagh, Dhaka, and partially relocated to the Savar Tannery Industrial Estate following the 2017 relocation mandate, has for decades discharged chromium-laden wastewater into the Buriganga and Shitalakshya rivers with demonstrably inadequate treatment infrastructure. A systematic survey of Buriganga river water (dry season, 2021--2023) recorded Cr(VI) concentrations spanning 0.01--3.54 mg/L (mean: 0.48 mg/L), with total dissolved chromium at industrial outfall points reaching 8.42 mg/L \cite{rahman2021,ahmed2022}. For context, the US EPA maximum contaminant level for total chromium in surface water is 0.1 mg/L; the WHO guideline for drinking water is 0.05 mg/L. In the Shitalakshya River (Narayanganj), which flows past Govt. Tolaram College, Cr(VI) concentrations range from 0.037 to 0.102 mg/L, with bioaccumulation in local fish tissue measured at 0.5--2 mg/kg dry weight \cite{islam2020,doem2022}. Dissolved oxygen (DO) in the Buriganga near tannery outfalls falls below 2 mg/L during the dry season, a level below the hypoxic threshold lethal to most fish species, with near-complete loss of aquatic macroinvertebrate diversity in affected zones. Communities dependent on these rivers for subsistence fishing, smallholder agriculture, and informal water access face compounding public health and economic consequences.

\subsection{Limitations of Existing Remediation Approaches}
Conventional Cr(VI) remediation methodologies each carry fundamental limitations that prevent their adoption as scalable, adaptive, and cost-effective solutions in low-resource settings. Chemical precipitation (reduction to Cr(III) using sodium bisulfite or ferrous sulfate at pH 2--3, followed by lime-induced Cr(OH)$_3$ precipitation) generates large volumes of hazardous chromium-hydroxide sludge requiring further processing, operates in narrow pH windows incompatible with river conditions, and exhibits no ability to modulate treatment intensity in response to fluctuating Cr(VI) concentrations \cite{sharma2017}. Electrocoagulation, the approach taken by a prior Bangladesh SJWP entry in 2022, requires continuous electrical input (30--50 W/m$^3$), is economically prohibitive for remote or rural deployment, and applies equal treatment intensity regardless of actual contamination levels \cite{heidmann2008}. Phytoremediation (hyperaccumulating plants such as \textit{Thlaspi caerulescens}) operates on timescales of months to years and is constrained by seasonal biomass availability \cite{ali2013}. Conventional unmodified bioremediation exploits endogenous chromate reductase activity in naturally occurring strains (e.g., \textit{Bacillus subtilis}, \textit{Lysinibacillus fusiformis}), but these organisms cannot report chromium presence, cannot be programmed to self-terminate after remediation, and lack containment guarantees for open-water environments \cite{viti2003}.

Most critically: \textit{none of these approaches operate as integrated sensing-acting-self-terminating systems}. They cannot distinguish between a heavily polluted sampling site and a recovered one, they cannot scale their response in proportion to pollution levels, and they provide no mechanism to prevent long-term biological persistence in the treatment environment.

\subsection{Synthetic Biology as a Remediation Paradigm}
Synthetic biology offers a categorically different framework. By encoding programmable genetic logic directly into a bacterial chassis, it is possible to engineer organisms that act not as passive chemical sponges but as autonomous bio-computers: sensing a specific signal (Cr(VI) concentration), computing a proportional response (NemA expression level), executing that response (enzymatic reduction), and self-terminating upon task completion (kill switch activation). This "sense-respond-terminate" architecture, in which all three modules are governed by the same regulatory protein (ChrB), is the conceptual core of the present design.

However, prior synthetic biology approaches to bioremediation have consistently underestimated or insufficiently quantified the risks associated with deploying Biologically Modified Organisms (BMOs) in open natural environments. Engineered bacteria are, by definition, living organisms subject to Darwinian selection. Any kill switch can fail through mutation; any auxotrophy can theoretically revert through horizontal gene transfer or compensatory evolution. A single escaped, reproductively competent BMO in the Buriganga River, carrying antibiotic resistance genes and a chromate-reductase operon, could represent an ecological and public health catastrophe of incalculable magnitude. \textbf{The first and most important design decision of this project is therefore absolute: we do not propose open-water deployment.} The platform operates exclusively within a physically sealed, closed-system microcosm (bioreactor vessel or sealed pond liner), and the computational biosafety modeling in Section 4 provides the mathematical argument that, under its modeling assumptions, even within this closed system the probability of a viable escape is well below the $10^{-15}$ acceptable-risk threshold.

\subsection{Novelty and Research Hypothesis}
The novelty of this work is not the individual components (ChrB biosensors \cite{branco2013}, NemA reductases \cite{williams2003}, and holin-endolysin kill switches \cite{chan2016} have each been independently described in the literature), but rather: \textbf{(i)} their integration into a single, three-plasmid, tri-modular system governed by a unified regulatory signal; \textbf{(ii)} the proposal, by analogy to OYE-family active-site mutagenesis precedent \cite{mowafy2010}, of a hypothetical NemA catalytic variant (NemA$^{*2+}$) with substantially improved kinetics, presented here as a hypothesis-by-analogy rather than a designed variant; and \textbf{(iii)} to our knowledge, the first application of Luria-Delbr\"uck fluctuation theory to rigorously quantify the evolutionary failure probability of a layered BMO containment architecture in a heavy-metal bioremediation BMO under Bangladesh-specific closed-system deployment parameters. We do not claim priority over Luria-Delbr\"uck fluctuation analysis in general synthetic-biology biocontainment, where related applications exist \cite{lee2012,stirling2017,rottinghaus2022}; the scoping claim is the specific application to heavy-metal bioremediation in Bangladesh tannery-effluent parameters.

\textbf{Central hypothesis:} A tri-modular synthetic biology platform integrating a ChrB-sfGFP biosensor, a NemA chromate reductase (including a hypothesized mutant variant), and a dual-trigger holin-endolysin kill switch, expressed in a double-auxotrophic \textit{E. coli} DH5$\alpha$ $\Delta$thyA $\Delta$dapA chassis within a closed-system microcosm, can, under its modeling assumptions, achieve: (1) $\geq$99\% Cr(VI) removal from 100 $\mu$M starting concentration within 48 hours; (2) $\geq$99\% programmed cell death within 72 hours of chromate depletion; and (3) a mathematical probability of containment breach below $10^{-15}$ over a 30-day closed-system deployment, as validated by ODE systems modeling, metabolic burden analysis, and Luria-Delbr\"uck evolutionary biosafety modeling with explicit sensitivity analysis of the mutation-rate input.

\section{Genetic Circuit Architecture and Plasmid Design}

\subsection{System Overview: Tri-Modular Design Philosophy}
The central engineering challenge of whole-cell bioremediation is achieving \textit{coherent signal integration across functional modules}. In many prior designs, sensing, effector expression, and kill switch elements are independently regulated, creating the risk of temporal desynchronization: for example, a kill switch activating prematurely before remediation is complete, or the biosensor becoming saturated and decoupling from NemA expression at high Cr(VI) concentrations.

Our design resolves this by making the ChrB regulatory protein the single master regulator of all three modules simultaneously. The ChrB protein, a transcriptional repressor from \textit{Pseudomonas aeruginosa} that binds the PchrB operator sequence in the apo (chromate-free) state and releases it upon Cr(VI) binding, governs:
\begin{itemize}
    \item \textbf{Module 1:} De-repression of PchrB $\rightarrow$ sfGFP (biosensor output)
    \item \textbf{Module 2:} De-repression of PchrB $\rightarrow$ NemA (remediation output)
    \item \textbf{Module 3:} Repression of the \textit{inverted} PchrB $\rightarrow$ holin/endolysin (kill switch \textit{suppression} during active chromate sensing)
\end{itemize}

This architecture ensures that (a) biosensing and remediation are always co-activated in proportion to Cr(VI) concentration; (b) the kill switch is actively suppressed \textit{only} during periods when Cr(VI) is above the PchrB threshold (i.e., when remediation is needed); and (c) the kill switch is released \textit{automatically} when chromate is depleted, the biochemical signal that treatment is complete. This is not a simple two-component system: it is a self-computing molecular AND gate.

\begin{table*}[!t]
\centering
\caption{Plasmid Design and Origins of Replication}
\begin{tabular}{@{}llllll@{}}
\toprule
\textbf{Plasmid} & \textbf{Module} & \textbf{Backbone} & \textbf{Origin} & \textbf{Marker} & \textbf{Copy No.} \\ \midrule
pChrB-sfGFP & Biosensor & pUC19 & pMB1 & AmpR & High ($\sim$500) \\
pNemA-His & Reductase & pET-28a & pBR322 & KanR & Med ($\sim$40) \\
pKillSwitch-DT & Kill Switch & pACYC184 & p15A & CmR & Low ($\sim$15) \\ \bottomrule
\end{tabular}
\end{table*}

The three plasmid backbones were specifically selected for mutual compatibility: pMB1 (pUC19) and pBR322 (pET-28a) are derived from the same ColE1 lineage but are incompatible with each other in high- versus medium-copy configurations; p15A (pACYC184) belongs to a completely orthogonal incompatibility group \cite{velappan2021}, allowing stable co-existence of all three plasmids in the same cell. The three antibiotic markers (AmpR, KanR, CmR) are also fully orthogonal. The copy number differential across the three plasmids is intentional and is what \textit{coordinately regulates} the relative per-cell expression of the three modules: the biosensor ($\sim$500 copies/cell on pUC19) produces a robust sfGFP signal even at low Cr(VI) concentrations; the reductase ($\sim$40 copies/cell on pET-28a) balances NemA production against metabolic burden; the kill switch ($\sim$15 copies/cell on pACYC184) prevents constitutive Holin leakage from overwhelming the CI434 repressor. We note explicitly that this is \textit{coordinate transcriptional regulation through shared promoter logic and chosen RBS strengths}, not strict stoichiometric coupling; the $\sim$12-fold pUC19:pET-28a copy-number gap is the intended design lever, and we have not yet demonstrated experimentally that it produces the predicted per-cell sfGFP:NemA ratio.

\subsection{Module 1: ChrB-sfGFP Chromium Biosensor}
\subsubsection{Molecular Mechanism of ChrB-Based Sensing}
The ChrB protein is a member of the LysR-type transcriptional regulator (LTTR) superfamily, originally characterized in \textit{Pseudomonas aeruginosa} PAO1 \cite{branco2013}. In the absence of chromate, the ChrB dimer binds the PchrB operator sequence (a 23-bp palindromic motif) upstream of the \textit{chrBCA} operon and sterically blocks RNA polymerase progression. Upon binding of chromate anion (CrO$_4^{2-}$) to the co-inducer binding domain of ChrB, the protein undergoes an allosteric conformational shift that dramatically reduces its operator affinity, releasing the PchrB promoter and permitting transcriptional read-through \cite{igem2013}. Published characterization data for the BBa\_K1149051 BioBrick part (ChrB + PchrB promoter) indicate a Hill coefficient of $n \approx 2.0$--$2.2$ for the dose-response relationship, reflecting the dimeric binding mechanism and cooperative co-inducer binding.

\begin{figure}[!t]
    \centering
    \includegraphics[width=\linewidth]{simulations/results/plasmid_module1.png}
    \caption{Module 1 Architecture: High-copy pUC19 backbone carrying the PchrB-sfGFP biosensor cassette.}
\end{figure}

The predicted limit of detection (LOD) is $\sim$80 nM Cr(VI). Fold-induction at 100 $\mu$M Cr(VI) is predicted at 8.2$\times$, and response time to 90\% of maximum signal is $\sim$3 hours post-exposure.

\subsection{Module 2: NemA Chromate Reductase}
\subsubsection{The NemA Enzyme: Structure and Mechanism}
NemA (N-ethylmaleimide reductase) is a native \textit{E. coli} NADH-dependent flavoprotein belonging to the Old Yellow Enzyme (OYE) superfamily \cite{williams2003}. It exhibits significant promiscuity towards chromate as an electron acceptor. The catalytic mechanism proceeds through a ping-pong bi-bi mechanism:
\begin{equation}
\text{Cr(VI)} + \text{NADH} + \text{H}^+ \xrightarrow{\text{NemA/FMN}} \text{Cr(III)} + \text{NAD}^+
\end{equation}

Published kinetic parameters for wild-type NemA on Cr(VI) substrate: $K_m$ = 48 $\mu$M, $k_{cat}$ = 0.39 s$^{-1}$ (1,404 h$^{-1}$), $k_{cat}/K_m$ = $1.1 \times 10^5$ M$^{-1}$s$^{-1}$ \cite{williams2003}.

\begin{figure}[!t]
    \centering
    \includegraphics[width=\linewidth]{simulations/results/plasmid_module2.png}
    \caption{Module 2 Architecture: Medium-copy pET-28a backbone carrying the PchrB-NemA remediation cassette.}
\end{figure}

\subsubsection{Hypothesized NemA$^{*2+}$ Variant (Hypothesis-by-Analogy)}
To further maximize remediation efficiency, we \textit{hypothesize} a NemA catalytic variant (NemA$^{*2+}$) based on structural analogy to characterized OYE-family engineering. The OYE superfamily active site features a conserved Tyr-His catalytic pair that anchors the FMN isoalloxazine ring and positions the substrate for hydride transfer. A second shell of hydrophobic residues (Trp, Phe, Ile) flanks the substrate-binding pocket and governs the geometry of substrate approach, directly influencing $K_m$. Mutagenesis studies in homologous OYE members (PETNR; morphinone reductase) have demonstrated that substituting these bulky second-shell residues with smaller side chains (e.g., Trp$\rightarrow$Ala, Phe$\rightarrow$Gly) expands the substrate-binding pocket and reduces steric hindrance for non-canonical electron acceptors like chromate, lowering $K_m$ without disrupting the core Tyr-His catalytic mechanism or $k_{cat}$ \cite{mowafy2010}.

\textbf{Important methodological caveat.} The specific NemA$^{*2+}$ kinetic parameters used in the comparative model ($K_m^* = 16\ \mu$M, $k_{cat}^* = 37.2\ \text{h}^{-1}$) are \textbf{assumed values}, extrapolated from the OYE-family active-site mutagenesis precedent reported in \cite{mowafy2010}. They are \textit{not} the output of any structural modeling pipeline (e.g., AlphaFold2 prediction, FoldX $\Delta\Delta G$ scan, Rosetta design, or molecular docking) executed in this study. No such pipeline is included in this repository. The hypothesis is presented as a plausible design target whose kinetic improvement over wild-type is consistent with the published OYE literature on second-shell pocket engineering, and is offered as motivation for future wet-lab or computational work that \textit{would} perform such modeling. The projected kinetic improvement, were such a mutant to be realized, would be a 3-fold reduction in $K_m$ and a 2-fold increase in $k_{cat}$ (in the operational units used by the circuit model), yielding a projected 6-fold improvement in $k_{cat}/K_m$ catalytic efficiency.

\subsection{Module 3: Dual-Trigger Holin-Endolysin Kill Switch}
Our design utilizes two independent triggering mechanisms combined in an AND-gate logic to trigger cell death via Holin-Endolysin \cite{chan2016}:
\begin{itemize}
    \item \textbf{Trigger 1 (Biochemical state sensor):} Chromate depletion de-represses the inverted PchrB promoter, activating holin-endolysin transcription.
    \item \textbf{Trigger 2 (Temporal timer):} A constitutive CI434-SsrA repressor provides an absolute time limit.
\end{itemize}

\begin{figure}[!t]
    \centering
    \includegraphics[width=\linewidth]{simulations/results/plasmid_module3.png}
    \caption{Module 3 Architecture: Low-copy pACYC184 backbone carrying the dual-trigger Holin-Endolysin kill switch.}
\end{figure}

The 24-hour protein half-life of CI434 is enforced by fusing the SsrA degradation tag to its C-terminus. By 72 hours post-inoculation, CI434-SsrA has decayed to $e^{-3\ln 2} = 12.5\%$ of its initial level, insufficient to maintain full repression of the kill-switch promoter. Combined with declining chromate (Trigger 1), this ensures Holin accumulation crosses the lysis threshold between 72--96 hours, regardless of residual chromium levels.

\subsection{The Chassis: \textit{E. coli} DH5$\alpha$ $\Delta$thyA $\Delta$dapA}
The chassis serves as an independent biocontainment layer via two full gene deletions. $\Delta$thyA requires exogenous thymine supplementation (50 $\mu$g/mL), and $\Delta$dapA requires diaminopimelic acid (DAP, 100 $\mu$M) for peptidoglycan cross-linking, preventing survival in natural freshwater, soil, or sediment environments. Reversion of either deletion by simple back-mutation is biologically impossible (there is no sequence to revert to); the only escape mode is horizontal gene transfer of an intact functional gene, addressed quantitatively in Section 4.

\section{Computational Methods}

All simulation code is implemented in Python 3.11 using NumPy, SciPy, and Matplotlib, and is available in the \texttt{simulations/} directory of the project repository. Every numerical constant used in any of the four pillar models is centralized in \texttt{simulations/parameters.py}, where each constant is tagged with its provenance: \texttt{source: [N]} (a specific reference from the manuscript bibliography), \texttt{ASSUMED} (an estimated/chosen value with an inline justification), or \texttt{PREDICTED} (an output of one of our own models, not an input). Reviewers and readers can audit the source/assumed status of every number in the body text by opening this single file.

\subsection{Pillar 1: ODE Systems Biology Model (\texttt{circuit\_ode\_model.py})}
The dynamic behavior of the tri-modular circuit was modeled as a coupled ODE system implemented in \texttt{circuit\_ode\_model.py} and integrated with \texttt{scipy.integrate.solve\_ivp} using the implicit \textbf{Radau} solver for stiff systems (relative tolerance $r_{tol} = 10^{-6}$, absolute tolerance $a_{tol} = 10^{-9}$). Six state variables were tracked over a 96-hour simulation window:

\begin{equation}
\mathbf{y} = [\text{Cr}_{ext},\ \text{NemA},\ \text{sfGFP},\ \text{CI434},\ \text{Holin},\ \text{Cells}]
\end{equation}

The right-hand side is given by the following coupled ODEs. \textbf{Chromate reduction} (NemA, Module 2):
\begin{equation}
\frac{d[\text{Cr}_{ext}]}{dt} = -\frac{k_{cat}^{\text{op}} \cdot [\text{NemA}] \cdot [\text{Cr}_{ext}]}{K_m^{\text{op}} + [\text{Cr}_{ext}]}
\end{equation}

\textbf{ChrB-driven protein production} (Modules 1 and 2):
\begin{equation}
\frac{d[\text{NemA}]}{dt} = k_{NemA} \cdot \frac{[\text{Cr}_{ext}]^n}{K_d^n + [\text{Cr}_{ext}]^n} - \delta_{NemA} \cdot [\text{NemA}]
\end{equation}
\begin{equation}
\frac{d[\text{sfGFP}]}{dt} = k_{GFP} \cdot \frac{[\text{Cr}_{ext}]^n}{K_d^n + [\text{Cr}_{ext}]^n} - \delta_{GFP} \cdot [\text{sfGFP}]
\end{equation}

\textbf{CI434-SsrA timer decay} (Module 3):
\begin{equation}
\frac{d[\text{CI434}]}{dt} = -\lambda_{CI434} \cdot [\text{CI434}], \quad \lambda_{CI434} = \frac{\ln 2}{t_{1/2}^{CI434}}
\end{equation}

\textbf{Kill switch Holin accumulation} (AND-gate logic, Module 3):
\begin{equation}
\frac{d[\text{Holin}]}{dt} = k_{Holin} \cdot \underbrace{\frac{K_d^n}{K_d^n + [\text{Cr}_{ext}]^n}}_{\text{Trigger 1: low Cr}} \cdot \underbrace{\frac{K_{CI}^2}{K_{CI}^2 + [\text{CI434}]^2}}_{\text{Trigger 2: low CI434}}
\end{equation}

\textbf{Cell population dynamics} (logistic growth + lysis):
\begin{equation}
\frac{d[\text{Cells}]}{dt} = \mu_{max} \cdot [\text{Cells}] \cdot \left(1 - \frac{[\text{Cells}]}{K_{carry}}\right) - k_{lysis} \cdot [\text{Cells}] \cdot \mathbb{1}_{[\text{Holin}] > \theta_{lysis}}
\end{equation}

\textbf{Important methodological note on the operational vs.\ published NemA $k_{cat}$.} The Michaelis-Menten reduction equation above uses an \textit{operational} per-cell rate constant $k_{cat}^{\text{op}} = 18.6\ \text{h}^{-1}$ that is approximately 75-fold lower than the published per-enzyme-molecule value ($k_{cat} = 0.39\ \text{s}^{-1} \approx 1404\ \text{h}^{-1}$) reported by Williams et al.\ 2003 \cite{williams2003}. The operational value is a reduced effective rate that, when combined with the \textit{also} operational cellular NemA concentration unit, reproduces the qualitative 96-hour dynamics of the tri-modular circuit. It is not a direct physical measurement and should not be compared on a per-enzyme basis to the published $k_{cat}$. The Pillar 2 comparative model (Section~\ref{sec:pillar2}) uses the same operational convention for wild-type and mutant, so the relative comparison there is meaningful on its own terms.

\subsection{Pillar 2: NemA vs.\ NemA$^{*2+}$ Kinetic Comparison (\texttt{nemA\_mutant\_kinetics.py})}
\label{sec:pillar2}
Wild-type NemA and the hypothesized NemA$^{*2+}$ variant were compared by integrating the substrate-only Michaelis-Menten ODE over a 24-hour window starting from $[\text{Cr(VI)}]_0 = 100\ \mu$M and with total enzyme $E_{total} = 5\ \mu$M. The Radau solver was used with the same tolerances as Pillar 1. Wild-type operational parameters: $k_{cat}^{WT} = 18.6\ \text{h}^{-1}$, $K_m^{WT} = 48\ \mu$M. Hypothesized NemA$^{*2+}$ parameters: $k_{cat}^{*} = 37.2\ \text{h}^{-1}$, $K_m^{*} = 16\ \mu$M. The mutant values are \textbf{assumed}, extrapolated from active-site pocket mutagenesis precedent in OYE-family enzymes \cite{mowafy2010} (see Section 2.3.3 for the full caveat on provenance). The script also computes and plots the velocity vs.\ substrate curves for both enzymes to illustrate the tail-end (low-substrate) efficiency gain that the lower $K_m$ confers.

\subsection{Pillar 3: Ribosome Allocation Metabolic Burden Model (\texttt{metabolic\_burden\_model.py})}
Cellular growth rate under circuit expression and Cr(VI) stress was computed using the ribosome-allocation framework of Scott et al.\ 2010 \cite{scott2010}. The instantaneous specific growth rate is:
\begin{equation}
\mu = \left(\phi_{max} - \phi_q - \phi_s\right) \cdot \kappa_t^{stressed}
\end{equation}
where $\phi_{max}$ is the maximum mass fraction of the proteome that can be ribosomes (0.55), $\phi_q$ is the housekeeping protein fraction (0.40, adjusted upward from the typical $\sim$0.27 to reflect the maintenance demands of a closed-pond / nutrient-poor environment), $\phi_s$ is the synthetic-circuit proteome fraction swept from 0 to 0.15, and $\kappa_t^{stressed}$ is the translational elongation efficiency under Cr(VI) stress. The stress model is a linear penalty: $\kappa_t^{stressed} = \kappa_t \cdot (1 - [\text{Cr}]/[\text{Cr}]_{max})$, with $\kappa_t = 3.5\ \text{h}^{-1}$ (a nutrient-poor pond value, lower than the $\sim$6--8 h$^{-1}$ of Scott et al.\ 2010's standard laboratory condition) and $[\text{Cr}]_{max} = 250\ \mu$M.

\subsection{Pillar 4: Luria-Delbr\"uck Evolutionary Biosafety Model (\texttt{biosafety\_mutation\_model.py})}
Evolutionary containment failure probability was computed using the Poisson approximation to Luria-Delbr\"uck fluctuation statistics. The probability of at least one escape mutant appearing in a population of $N$ total cells over $G$ generations, given per-cell-per-generation failure probability $p$:
\begin{equation}
P_{\text{escape}} = 1 - e^{-N \cdot G \cdot p}
\end{equation}

For the combined kill-switch + double-auxotrophy containment layers, the per-division failure probability is the product of three independent failure modes (a cell must simultaneously fail all three to escape):
\begin{equation}
p_{\text{combined}} = p_{KS\,fail} \times p_{\Delta thyA\,HGT} \times p_{\Delta dapA\,HGT}
\end{equation}
where $p_{KS\,fail} = \mu_{bp} \cdot L_{target}$ (per-bp mutation rate $\times$ kill-switch target size), and $p_{\Delta thyA\,HGT}$, $p_{\Delta dapA\,HGT}$ are horizontal gene transfer probabilities for intact functional \textit{thyA} and \textit{dapA} genes from environmental donors. The script produces two output figures: a 365-day long-run view and a 30-day deployment-window view cropped to the deployment window that the hypothesis explicitly targets. The latter is the figure referenced in Section 4.

\subsection{Pillar 4b: Biosafety Sensitivity Analysis (\texttt{biosafety\_sensitivity.py})}
The strongest claim in the manuscript (that $P_{\text{escape}}$ remains below an acceptable-risk threshold of $10^{-15}$) is sensitive to the assumed spontaneous mutation rate $\mu_{bp}$, which the literature places anywhere in the range $10^{-10}$ to $10^{-9}$ per bp per generation. \texttt{biosafety\_sensitivity.py} sweeps $\mu_{bp}$ over scaling factors $\times$0.01, $\times$0.1, $\times$1 (baseline), $\times$10, $\times$100, holding the deployment scale fixed at the 1,000 L / $10^{12}$ cell baseline, and reports $P_{\text{escape}}$ at the 30-day window for each scaling. It also sweeps reactor volume over 10 L, 100 L, 1,000 L, and 10,000 L at the baseline mutation rate. The results are reported in Section~\ref{sec:sens}.

\subsection{Parameter Table}
\begin{table*}[!t]
\centering
\caption{Centralized parameter set for all four computational pillars. Full provenance and per-constant inline justification are in \texttt{simulations/parameters.py}.}
\label{tab:params}
\small
\begin{tabular}{@{}lllll@{}}
\toprule
\textbf{Symbol} & \textbf{Value} & \textbf{Units} & \textbf{Provenance} & \textbf{Used by} \\ \midrule
$K_d$ (PchrB) & 0.1 (100 nM) & $\mu$M & source: \cite{igem2013} & Pillar 1 \\
$n$ (Hill coeff.) & 2 & --- & source: \cite{igem2013} & Pillar 1 \\
$K_m^{WT}$ (operational) & 48 & $\mu$M & source: \cite{williams2003} & Pillars 1, 2 \\
$k_{cat}^{WT}$ (operational) & 18.6 & h$^{-1}$ & ASSUMED (operational rate) & Pillars 1, 2 \\
$k_{cat}^{WT}$ (per-enzyme) & 0.39 / 1404 & s$^{-1}$ / h$^{-1}$ & source: \cite{williams2003} & Pillar 2 (ref.) \\
$K_m^{*}$ (NemA$^{*2+}$) & 16 & $\mu$M & ASSUMED (extrapolated from \cite{mowafy2010}) & Pillar 2 \\
$k_{cat}^{*}$ (NemA$^{*2+}$) & 37.2 & h$^{-1}$ & ASSUMED (extrapolated from \cite{mowafy2010}) & Pillars 1, 2 \\
$t_{1/2}^{CI434}$ & 24 & h & source: SsrA degron literature & Pillar 1 \\
$\phi_{max}$ & 0.55 & --- & source: \cite{scott2010} & Pillar 3 \\
$\phi_q$ & 0.40 & --- & ASSUMED (nutrient-poor pond) & Pillar 3 \\
$\kappa_t$ & 3.5 & h$^{-1}$ & ASSUMED (nutrient-poor pond) & Pillar 3 \\
$[\text{Cr}]_{max}$ & 250 & $\mu$M & ASSUMED & Pillar 3 \\
$\mu_{bp}$ & $10^{-9}$ & per bp per gen & ASSUMED (Drake \cite{drake1991} range) & Pillar 4 \\
$L_{target}$ (kill switch) & 500 & bp & ASSUMED (per \cite{chan2016} arch.) & Pillar 4 \\
$p_{\Delta thyA\,HGT}$ & $10^{-12}$ & per division & ASSUMED (no single primary citation) & Pillar 4 \\
$p_{\Delta dapA\,HGT}$ & $10^{-12}$ & per division & ASSUMED (no single primary citation) & Pillar 4 \\
$V$ (deployment volume) & 1000 & L & source: Sec.~4 & Pillar 4 \\
Max cell density & $10^9$ & cells/L & source: Sec.~4 & Pillar 4 \\
Generations/day & 4 & gen/day & ASSUMED (nutrient-poor pond) & Pillar 4 \\
Deployment window & 30 & days & source: Sec.~1.4 (hypothesis) & Pillar 4 \\
De minimis risk threshold & $10^{-15}$ & --- & source: Sec.~4 & Pillar 4 \\
\bottomrule
\end{tabular}
\end{table*}

\FloatBarrier

\section{Results}

\subsection{Circuit Systems Dynamics: 96-Hour ODE Simulation}

\begin{figure*}[!t]
    \centering
    \includegraphics[width=0.8\textwidth]{simulations/results/integrated_96h_simulation.png}
    \caption{96-hour ODE simulation of tri-modular circuit dynamics. (Top-left) Cr(VI) depletion by NemA-catalyzed reduction; (top-right) sfGFP biosensor fluorescence as a function of time; (bottom-left) CI434-SsrA timer decay and Holin accumulation crossing the lysis threshold $\theta_{lysis}=5$ (red dashed line); (bottom-right) live-cell density on a logarithmic scale. Generated by \texttt{circuit\_ode\_model.py}.}
    \label{fig:ode}
\end{figure*}

The 96-hour coupled-ODE simulation (Fig.~\ref{fig:ode}) reveals precise temporal coordination across all three modules. Starting from an initial condition of 100 $\mu$M Cr(VI) and $1 \times 10^7$ CFU/mL, the system executes the remediation protocol in three distinct phases. \textbf{Phase 1 (0--8 h):} sfGFP fluorescence rises rapidly as ChrB releases the PchrB operator; the kill switch is held quiescent. \textbf{Phase 2 (8--44 h):} NemA reaches steady-state concentration; Cr(VI) decreases approximately linearly at $\sim$2.3 $\mu$M/h, falling below 100 nM by $t = 44$ h. \textbf{Phase 3 (44--96 h):} The kill-switch AND-gate fires; Holin crosses $\theta_{lysis} = 5$ at $t \approx 72$ h, and live-cell density drops by $>$4 orders of magnitude by 96 h.

\subsection{NemA vs.\ NemA$^{*2+}$ Kinetic Comparison (Hypothesized)}

\begin{figure*}[!t]
    \centering
    \includegraphics[width=0.8\textwidth]{simulations/results/nemA_mutant_kinetics.png}
    \caption{Comparative Michaelis-Menten kinetics of wild-type NemA and the hypothesized NemA$^{*2+}$ variant (parameters assumed by analogy to OYE-family active-site mutagenesis precedent \cite{mowafy2010}; see Section 2.3.3). (Left) Time-course depletion of 100 $\mu$M Cr(VI); (right) initial reaction velocity as a function of substrate concentration, showing the $\sim$3-fold decrease in apparent $K_m$. Generated by \texttt{nemA\_mutant\_kinetics.py}.}
    \label{fig:kinetics}
\end{figure*}

The kinetic comparison between wild-type NemA and the hypothesized NemA$^{*2+}$ variant (parameters assumed by analogy to OYE-family active-site mutagenesis precedent \cite{mowafy2010}; see Section 2.3.3) shows the projected performance difference (Fig.~\ref{fig:kinetics}). Under the model's assumptions, wild-type NemA reduces 100 $\mu$M Cr(VI) to below the US EPA surface water limit (0.1 mg/L $\approx$ 1.9 $\mu$M) in approximately 43 h. The hypothesized NemA$^{*2+}$ variant, with its three-fold lower $K_m$ and doubled $k_{cat}$ (in the operational units used by the circuit model), achieves the same endpoint in $\sim$11 h, a four-fold reduction in \textit{predicted} treatment time. This 4-fold figure is \textbf{conditional on the caveat} in Section 2.3.3: the kinetic parameters are assumed, not derived from a structural modeling pipeline executed in this study.

At environmentally relevant Cr(VI) concentrations (1--10 $\mu$M, characteristic of moderately polluted reaches of the Buriganga and Shitalakshya), the hypothesized mutant operates at $\sim$40--75\% of $V_{max}$, whereas the wild-type operates at only $\sim$2--17\%. This differential becomes decisive in low-concentration field scenarios where most of the reaction time is spent in the $S \ll K_m$ regime.

\subsection{Metabolic Feasibility: Circuit Burden Analysis}

\begin{figure*}[!t]
    \centering
    \includegraphics[width=0.8\textwidth]{simulations/results/metabolic_burden_model.png}
    \caption{Ribosome-allocation metabolic burden model. Growth rate $\mu$ (h$^{-1}$) is plotted against synthetic-circuit expression level (\% of total proteome) at four environmental Cr(VI) concentrations. The dashed vertical line marks the predicted 4.5\% burden of our circuit at full induction. Generated by \texttt{metabolic\_burden\_model.py}.}
    \label{fig:burden}
\end{figure*}

The ribosome-allocation metabolic model (Fig.~\ref{fig:burden}) demonstrates that, under its assumptions, the proposed circuit is metabolically sustainable under all relevant environmental chromium concentrations. At $\phi_S = 4.5\%$, growth rate is reduced by approximately 33\% relative to the unstressed, unburdened baseline (0.5 h$^{-1}$ $\to$ 0.34 h$^{-1}$ at 0 $\mu$M Cr(VI)). Growth remains positive under all four stress conditions tested (0, 20, 50, 100 $\mu$M Cr(VI)), indicating that the chassis continues to replicate and remediate.

\subsection{Evolutionary Biosafety: Luria-Delbr\"uck Containment Analysis}

\begin{figure*}[!t]
    \centering
    \includegraphics[width=0.8\textwidth]{simulations/results/biosafety_mutation_30day.png}
    \caption{Luria-Delbr\"uck evolutionary biosafety analysis over the 30-day deployment window. Probability of containment failure is plotted for two scenarios: kill-switch-only containment (red dashed) and the full 5-layer containment including double auxotrophy (solid green). The orange dotted line marks the $10^{-15}$ `de minimis' acceptable-risk threshold. The full containment curve sits at $P_{\text{escape}} \approx 6 \times 10^{-17}$ (see Eq.~12 and surrounding text for the IEEE-754 arithmetic note). A 365-day long-run supplementary view is provided in Fig.~\ref{fig:biosafety-long} below. Generated by \texttt{biosafety\_mutation\_model.py}.}
    \label{fig:biosafety}
\end{figure*}

\begin{figure*}[!t]
    \centering
    \includegraphics[width=0.8\textwidth]{simulations/results/biosafety_mutation_model.png}
    \caption{Luria-Delbr\"uck evolutionary biosafety analysis --- 365-day long-run supplementary view. Same model as Fig.~\ref{fig:biosafety}, plotted out to 365 days to show the kill-switch-only curve crossing the $10^{-15}$ `de minimis' threshold at approximately day 80. The full 5-layer containment curve remains far below the threshold over the full year. This long-run view is supplementary; the headline 30-day result is reported in Fig.~\ref{fig:biosafety} and Eq.~12.}
    \label{fig:biosafety-long}
\end{figure*}

The Luria-Delbr\"uck biosafety analysis (Fig.~\ref{fig:biosafety}) provides the quantitative argument for the layered containment architecture. We model a 1,000-liter closed pond at carrying capacity ($1 \times 10^9$ cells/L, total $1 \times 10^{12}$ cells), with 4 generations/day and 30 days of operation ($1.2 \times 10^{14}$ total cell divisions). Per-division failure probabilities:
\begin{itemize}
    \item \textbf{Kill-switch loss-of-function:} $P_{ks} = \mu_{bp} \times L_{target} = 10^{-9} \times 500 = 5 \times 10^{-7}$ per division, where $\mu_{bp} = 10^{-9}$ per bp per generation and $L_{target} = 500$ bp.
    \item \textbf{Auxotrophy reversion (HGT):} $P_{rev} = 10^{-12}$ per division per locus (theoretical lower bound for HGT-driven reconstitution).
    \item \textbf{Combined per-division escape:} $P_{\text{escape/gen}} = 5 \times 10^{-7} \times 10^{-12} \times 10^{-12} = 5 \times 10^{-31}$.
\end{itemize}

The 30-day integrated probability:
\begin{equation}
P_{\text{escape}}^{30d} = 1 - e^{-N_{total} \cdot P_{\text{escape/gen}}} = 1 - e^{-1.2 \times 10^{14} \times 5 \times 10^{-31}} = 1 - e^{-6 \times 10^{-17}} \approx 6 \times 10^{-17}
\end{equation}

\textbf{Numerical implementation note.} For arguments this small, evaluating $1 - e^{-x}$ directly in IEEE-754 double precision suffers from catastrophic cancellation (the $\exp(-6 \times 10^{-17})$ result is dominated by roundoff in the leading `1', and a naive Python/NumPy evaluation returns $\approx 1.11 \times 10^{-16}$ rather than the analytically correct $6 \times 10^{-17}$). The simulation scripts (\texttt{biosafety\_mutation\_model.py}, \texttt{biosafety\_sensitivity.py}) therefore use the first-order Taylor expansion $1 - e^{-x} \approx x$ for $x < 10^{-10}$, which agrees with the algebra shown in Eq.~12 to all displayed significant figures. This is a numerical-analysis detail, not a physics choice, and would be invisible in any practical deployment-safety argument.

This is approximately \textbf{17 orders of magnitude} below the $10^{-15}$ `de minimis' threshold. The kill-switch-only curve, when extrapolated to the long-run view in Fig.~\ref{fig:biosafety-long} (\texttt{results/biosafety\_mutation\_model.png}, the 365-day supplementary plot), crosses the threshold at approximately day 80; this is not visible in the cropped 30-day Fig.~\ref{fig:biosafety} shown here, but the relevant long-run data is included in the supplementary figure for completeness. The double auxotrophy thus provides an essential second-tier failsafe.

\subsubsection{Sensitivity of the Containment Claim to Mutation-Rate and Deployment-Scale Assumptions}
\label{sec:sens}
The headline $P_{\text{escape}} < 10^{-15}$ claim depends on the specific value of the per-bp spontaneous mutation rate ($\mu_{bp} = 10^{-9}$), which the literature places in a range spanning $10^{-10}$ to $10^{-9}$ \cite{drake1991,lee2012}. To test the robustness of the claim, we re-ran the biosafety calculation with $\mu_{bp}$ scaled by $\times$0.01, $\times$0.1, $\times$1 (baseline), $\times$10, and $\times$100, and swept reactor volume across 10 L, 100 L, 1,000 L (baseline), and 10,000 L. Results are in Fig.~\ref{fig:sens} and \texttt{results/sensitivity\_analysis.csv}.

\begin{figure*}[!t]
    \centering
    \includegraphics[width=0.8\textwidth]{simulations/results/sensitivity_analysis.png}
    \caption{Sensitivity of the 30-day $P_{\text{escape}}$ to (left) a 0.01$\times$--100$\times$ scaling of the assumed per-bp mutation rate, and (right) a 10 L--10,000 L sweep of the deployment reactor volume. The $10^{-15}$ acceptable-risk threshold is shown in orange; $P_{\text{escape}} = 1$ (certain failure) in red. Generated by \texttt{biosafety\_sensitivity.py}.}
    \label{fig:sens}
\end{figure*}

\textbf{Key findings (sensitivity sweep, 30-day window):}
\begin{itemize}
    \item At baseline ($\mu_{bp} = 10^{-9}$): $P_{\text{escape}} = 6.00 \times 10^{-17}$.
    \item At $\times$0.1 mutation rate: $P_{\text{escape}} = 6.00 \times 10^{-18}$.
    \item At $\times$10 mutation rate: $P_{\text{escape}} = 6.00 \times 10^{-16}$, still below $10^{-15}$ by $\sim$1.7$\times$.
    \item \textbf{At $\times$100 mutation rate: $P_{\text{escape}} = 6.00 \times 10^{-15}$, exactly 6$\times$ above the $10^{-15}$ threshold.}
    \item Across the reactor volume sweep (10 L--10,000 L) at baseline mutation rate, $P_{\text{escape}}$ remains below $10^{-15}$ for all four volumes.
\end{itemize}

\textbf{Interpretation.} The headline $P_{\text{escape}} < 10^{-15}$ claim is robust to an order-of-magnitude error in the mutation-rate assumption and to a 10$\times$ scale-up of the deployment reactor. It is \textbf{not} robust to a two-order-of-magnitude error in $\mu_{bp}$: under that scenario the model places $P_{\text{escape}}$ at $6.00 \times 10^{-15}$, which is \textit{above} the $10^{-15}$ threshold and is therefore, by the model's own logic, a containment failure. This is not a finding to be smoothed over --- it is the \textit{useful} finding, because it shows the analysis was not tuned to produce reassuring numbers. The sensitivity analysis directly motivates the most important experimental follow-up: \textbf{the headline containment claim is contingent on a direct measurement of the actual per-bp mutation rate in the deployed strain under the deployed growth conditions} (closed-pond, burdened, with the tri-modular plasmid set in place). Until that measurement exists, the $10^{-15}$ risk threshold should be read as the boundary the model sits at under its current assumption, not as a margin of safety; a future in-vitro mutation-rate assay (e.g., a Luria-Delbr\"uck fluctuation test on the deployed strain, or a mutation-accumulation whole-genome sequencing line) is the single experiment that would convert the headline from a model output to a validated claim. See Section 5.3 (Limitation 5) for the related HGT-probability caveat.

\FloatBarrier

\section{Discussion}

\subsection{The Engineering Elegance of a Unified ChrB Regulatory Signal}
The central design philosophy of this system (making ChrB the single master regulator of all three modules) deserves explicit discussion, as it is the feature that distinguishes this platform from patchwork multi-component approaches. By coupling biosensing, remediation, and kill switch suppression to the same protein-DNA interaction, the system achieves inherent temporal synchronization:
\begin{itemize}
    \item \textbf{Remediation is active if and only if the biosensor is active.}
    \item \textbf{Kill switch suppression is active if and only if the biosensor is active.}
    \item \textbf{The kill switch fires when, and only when, the remediation task is complete.}
\end{itemize}
This self-computing architecture minimizes the number of independent regulatory elements that could independently fail, and creates a system where failure in any one module has compensatory consequences in another.

\subsection{Why Closed-System Deployment Is Non-Negotiable}
The Luria-Delbr\"uck analysis provides quantitative confirmation of what biosafety intuition already suggests: that releasing a BMO into the Buriganga River is irresponsible. In an open-river system, the bacterium would encounter (a) unlimited genetic donors for HGT, (b) unlimited space and nutrients for population expansion, and (c) no physical containment boundary. The combined effect is that $N_{total} \times G$ in an open water is effectively unbounded, and any finite $p_{\text{combined}}$, however small, yields $P_{\text{escape}} \to 1$ over sufficient time.

\subsection{Limitations}
This study has several important limitations.

(1) \textbf{Purely computational, no wet-lab validation.} No experimental or wet-lab data were generated in this project; all results are derived from the simulation models in \texttt{simulations/}. Physical validation in a BSL-1 facility is required before any of the engineering claims can be defended as physical performance.

(2) \textbf{NemA$^{*2+}$ is a hypothesis, not a validated variant.} The kinetic parameters used for the NemA$^{*2+}$ variant ($K_m^* = 16\ \mu$M, $k_{cat}^* = 37.2\ \text{h}^{-1}$) are \textbf{assumed}, extrapolated by analogy to active-site pocket mutagenesis precedent in the OYE family \cite{mowafy2010}. No structural modeling pipeline was performed in this study to derive these values; no such pipeline is included in this repository. The hypothetical 4-fold treatment-time reduction reported in Section 4.2 is therefore conditional on the assumption that the OYE active-site engineering precedent transfers to NemA.

(3) \textbf{Real effluent has co-contaminants not modeled.} Bangladesh tannery effluent contains a complex mixture of metals (Cr, Pb, Cd, As, Cu, Zn), high salinity, sulfate, sulfide, ammonia, dissolved organic loads, and variable pH. Our model treats Cr(VI) as the sole toxicant in an otherwise benign medium. The actual toxicity, growth-rate, and biosensor response in real effluent will differ from the model prediction.

(4) \textbf{Cr(III) downstream fate not addressed.} This work models Cr(VI) $\rightarrow$ Cr(III) reduction but does not address the subsequent fate of the Cr(III) product. At neutral-to-alkaline pH and in the presence of oxidants, Cr(III) can re-oxidize back to Cr(VI). A complete bioremediation design would need to address Cr(III) stabilization or sequestration.

(5) \textbf{Mutation-rate assumption is the key uncertainty in the containment claim.} As shown in the sensitivity analysis (Section 4), the headline $P_{\text{escape}} < 10^{-15}$ result is robust to a 10$\times$ error in the per-bp mutation rate but \textbf{not} to a 100$\times$ error. The literature on spontaneous \textit{E. coli} mutation rate is itself only constrained to within roughly one order of magnitude across studies, and the specific value for our chassis under our (closed-pond, burdened) growth conditions is not directly measured. The HGT reversion probabilities for $\Delta$thyA and $\Delta$dapA are order-of-magnitude estimates with no single primary citation. The $10^{-15}$ risk threshold should be read as a model output at the assumed input parameters, not a deep safety margin.

(6) \textbf{Coordinate vs.\ stoichiometric regulation.} The three-plasmid design uses copy-number differentials (pUC19 $\sim$500, pET-28a $\sim$40, pACYC184 $\sim$15) as the \textit{coordinate} transcriptional lever for the three modules. We have not experimentally verified that the resulting per-cell sfGFP:NemA:Holin ratios match the model assumptions.

\section{Conclusions}

We present a fully computational, multi-methodology biosafety-validated design of a tri-modular genetic circuit for Cr(VI) bioremediation in closed-system microcosms, with all parameters, assumptions, and provenance transparently documented in \texttt{simulations/parameters.py}. The key contributions are:
\begin{enumerate}
    \item A unified ChrB regulatory architecture that \textit{coordinately} regulates biosensing, enzymatic remediation, and kill switch suppression through a single regulatory protein.
    \item A hypothesis-by-analogy for a NemA$^{*2+}$ catalytic variant that, if realized, would offer a 6-fold improvement in $k_{cat}/K_m$ and reduce predicted treatment time from $\sim$43 to $\sim$11 hours. The kinetic parameters are \textit{assumed}, extrapolated from OYE-family active-site mutagenesis precedent \cite{mowafy2010} rather than derived from a structural modeling pipeline; the manuscript does not claim NemA$^{*2+}$ as a designed variant.
    \item A ribosome-allocation metabolic model showing that, under its modeling assumptions, the $\sim$4.5\% proteome burden of the synthetic circuit does not impair cell viability during the remediation window.
    \item A Luria-Delbr\"uck evolutionary biosafety analysis that, under the assumed input parameters, places the 30-day $P_{\text{escape}}$ of a viable escapee in a 1,000-liter closed system at $6 \times 10^{-17}$, well below the $10^{-15}$ acceptable-risk threshold.
    \item A Pillar 4b sensitivity analysis demonstrating that the headline $P_{\text{escape}} < 10^{-15}$ claim is robust to a 10$\times$ error in the per-bp mutation rate assumption and to a 10$\times$ scale-up of the deployment reactor, but is \textbf{not} robust to a 100$\times$ error in the mutation rate. The mutation rate is the single most consequential uncertainty in the containment claim, and the $10^{-15}$ threshold is a soft target rather than a deep safety margin.
\end{enumerate}

This computational framework does not merely describe what the system should do --- it provides a transparent mathematical foundation to evaluate both the design's expected performance and the sensitivity of its safety claim to its key input assumptions, and to identify the specific experiments (mutation-rate measurement under deployment conditions; AlphaFold2 / FoldX structural modeling of NemA$^{*2+}$; BSL-1 chassis co-transformation) that would convert it from a design study into a validated one. All such experimental validation is explicitly out of scope for this manuscript.

\section{Code and Data Availability}
All simulation code is contained in the public repository accompanying this manuscript. No experimental or wet-lab data were generated in this project --- all results are derived from the computational models in \texttt{simulations/}. The full set of numerical constants used by the models, with their provenance (literature-sourced vs.\ assumed/estimated), is centralized in \texttt{simulations/parameters.py}. This file is the single source of truth for every number in the project; each constant is tagged inline as \texttt{source: [N]}, \texttt{ASSUMED} (with justification), or \texttt{PREDICTED} (an output of one of our own models). The repository is archived at Zenodo: \textit{[DOI to be inserted after first Zenodo release]}.

\begin{thebibliography}{99}
\bibitem{iarc1990} IARC Working Group. \textit{Chromium, Nickel and Welding}. IARC Monographs Vol. 49. (1990).
\bibitem{eastmond2008} Eastmond DA, et al. \textit{Mutat Res.} 650(2):107--122 (2008).
\bibitem{rahman2021} Rahman MM, et al. \textit{Environ Monit Assess.} 193(7):420 (2021).
\bibitem{ahmed2022} Ahmed MF, et al. \textit{J Environ Sci.} 115:112--123 (2022).
\bibitem{islam2020} Islam MS, et al. \textit{J Hazard Mater.} 399:123074 (2020).
\bibitem{doem2022} DOEM. \textit{Annual Water Quality Monitoring Report: Narayanganj District}. (2022).
\bibitem{sharma2017} Sharma SK, et al. \textit{Desalin Water Treat.} 59:1--12 (2017).
\bibitem{heidmann2008} Heidmann I, Calmano W. \textit{J Hazard Mater.} 152(3):934--941 (2008).
\bibitem{ali2013} Ali H, et al. \textit{Chemosphere.} 91(7):869--881 (2013).
\bibitem{viti2003} Viti C, et al. \textit{Environ Microbiol.} 5(5):348--356 (2003).
\bibitem{branco2013} Branco R, et al. \textit{J Bacteriol.} 195(10):2054--2065 (2013).
\bibitem{williams2003} Williams RE, et al. \textit{Biochem Biophys Res Commun.} 290(1):249--254 (2003).
\bibitem{chan2016} Chan CTY, et al. \textit{Nat Chem Biol.} 12(2):82--86 (2016).
\bibitem{igem2013} iGEM Edinburgh 2013 Team. \textit{iGEM Parts Registry}. BBa\_K1149051. (2013).
\bibitem{mowafy2010} Mowafy AM, et al. \textit{Protein Sci.} 19(7):1283--1296 (2010).
\bibitem{lee2012} Lee JW, et al. Absolute quantitative analysis of the \textit{E. coli} spontaneous mutation rate per generation, per nucleotide, per cell. \textit{PNAS.} 109(41):E2774--E2783 (2012).
\bibitem{velappan2021} Velappan N, et al. Plasmid incompatibility: a comprehensive review. \textit{Plasmid.} 114:102557 (2021).
\bibitem{scott2010} Scott M, Gunderson CW, Mateescu EM, Zhang Z, Hwa T. Interdependence of cell growth and gene expression: origins and consequences. \textit{Science.} 330(6007):1099--1102 (2010).
\bibitem{stirling2017} Stirling F, et al. Rational design of evolutionarily stable microbial kill switches. \textit{Mol Cell.} 68(4):686--697 (2017).
\bibitem{rottinghaus2022} Rottinghaus AG, et al. Computational design of a biocontainment system for engineered microorganisms. \textit{ACS Synth Biol.} 11(1):134--145 (2022).
\bibitem{drake1991} Drake JW. A constant rate of spontaneous mutation in DNA-based microbes. \textit{Proc Natl Acad Sci USA.} 88(16):7160--7164 (1991).
\end{thebibliography}

\end{document}
